\documentclass[11pt]{article}
\usepackage[T1]{fontenc}
\usepackage[utf8]{inputenc}
\usepackage[margin=26mm]{geometry}
\usepackage{amsmath,amssymb,amsthm,booktabs,array,longtable}
\usepackage[colorlinks=true,allcolors=blue,breaklinks=true]{hyperref}
\newtheorem{theorem}{Theorem}
\newtheorem{lemma}[theorem]{Lemma}
\newtheorem{corollary}[theorem]{Corollary}
\newcommand{\C}{\mathbb C}
\newcommand{\R}{\mathbb R}
\newcommand{\tr}{\operatorname{tr}}
\newcommand{\diag}{\operatorname{diag}}
\title{Answering the finite-outcome question of Aloy et al.\ (2024),\\
arXiv:2312.09278v2, end of \S IV.B.3, p.~17:\\
an extension of their binary spin-one theorem}
\author{Haobo Ma\\{\small ChronoAI Pte Ltd, Singapore}
\and Wenlin Zhang\\{\small National University of Singapore}\\
{\small ChronoAI Pte Ltd, Singapore}}
\date{November 11, 2026}
\begin{document}
\maketitle
\begin{abstract}
Extending the binary-outcome theorem of Aloy, Galley, Jones, Ludescher and
M\"uller (2024), we answer yes to their unnumbered question at the end of
\S IV.B.3 of arXiv:2312.09278v2, p.~17 (\S IV.2.3 in the arXiv HTML):
$\mathcal Q_1^{\mathcal A}=\mathcal R_1^{\mathcal A}$ for every finite
outcome set with at least three elements. A positive Gram matrix for each
degree-two probability polynomial is averaged with its reversed transpose.
This preserves each polynomial and makes the sum of the matrices diagonal.
Diagonal normalization then supplies one pure qutrit state and one POVM
for the whole family, including singular boundary cases.
\end{abstract}
\noindent\textbf{License.} CC BY 4.0 (\texttt{cc-by-4.0}).

\section{The question and the result}
Let $\mathcal A$ be a finite nonempty set. Following Definitions~1 and~2 of
\cite{AGJLM}, write $\mathcal R_1^{\mathcal A}$ for the families
$p=(p_a)_{a\in\mathcal A}$ of real trigonometric polynomials of degree at
most two such that
\begin{equation}\label{probability}
 p_a(\theta)\geq0,\qquad \sum_{a\in\mathcal A}p_a(\theta)=1
 \quad(\theta\in\R).
\end{equation}
Each $p_a\leq1$ follows from these conditions.
The set $\mathcal Q_1^{\mathcal A}$ consists of quantum realizations
\begin{equation}\label{born}
 p_a(\theta)=\tr\bigl(U_\theta\rho U_\theta^*E_a\bigr),\qquad
 U_\theta=\bigoplus_{j=-1}^{1}e^{ij\theta}I_{n_j},
\end{equation}
where $n_j$ are finite nonnegative integers, $\rho\geq0$ has trace one,
and $E_a\geq0$, $\sum_a E_a=I$. The state and all effects are fixed as
$\theta$ varies. The source says ``do not currently know whether'' the equality
holds for $|\mathcal A|\geq3$ at the locator above and repeats the question
in \S VII, p.~27. Its Theorem~6 proves the binary equality, already present
in the first preprint of December 14, 2023. Its Theorem~1
already shows that quantum correlations admit pure states in dimension
$2J+1$; the issue here is whether every general spin-one family is quantum.
Definitions~1 and~2 in that source assume at least two outcomes. We extend
the same definitions to a one-element set, for which the only family is $p=1$.

\begin{theorem}\label{main}
For every finite nonempty $\mathcal A$,
$\mathcal Q_1^{\mathcal A}=\mathcal R_1^{\mathcal A}$.
Every family in this set has a realization on $\C^3$ with
$U_\theta=\diag(e^{-i\theta},1,e^{i\theta})$ and a pure state.
The state can be chosen as $\rho=\psi\psi^*$, where
\[
 \psi=(\sqrt t,\sqrt{1-2t},\sqrt t)^T,
 \qquad 0\leq t\leq\tfrac12.
\]
The parameter $t$ and the POVM may depend on the family $p$.
\end{theorem}

We give all steps of the construction. In particular, scalar positivity
alone is not the missing normalization step: the individual Gram matrices
must have a diagonal sum. The diagonal-sum constraint also appears in the
semidefinite characterization of fixed-spectrum quantum timelines in
\cite[\S III.B, Eqs.~(44)--(47)]{Navascues}. Their result characterizes
quantum timelines by positive matrices with diagonal sum; it does not identify
all normalized nonnegative degree-two families with those timelines.
Here the required diagonal sum is obtained for every family
satisfying \eqref{probability}.

\section{Positive Gram matrices and simultaneous normalization}
Put $v(\theta)=(e^{-i\theta},1,e^{i\theta})^T$ and
$q_Q(\theta)=v(\theta)^*Qv(\theta)$ for a $3\times3$ matrix $Q$.

\begin{lemma}\label{scalar}
Every nonnegative real trigonometric polynomial $f$ of degree at most two
has a representation $f=q_Q$ with $Q\geq0$.
\end{lemma}
This is the degree-two case of the scalar factorization consequence in
\cite[Lemma~9]{AGJLM}; the following elementary proof also fixes the Gram
convention used below.
\begin{proof}
Choose a minimum $m=f(\alpha)\geq0$, which exists by continuity and
periodicity. Write
\[
 g(s)=f(s+\alpha)-m=C_0+C_1\cos s+S_1\sin s
                       +C_2\cos(2s)+S_2\sin(2s).
\]
Since $g(0)=g'(0)=0$, we have $C_1=-C_0-C_2$ and $S_1=-2S_2$.
The double-angle identities therefore give
\begin{equation}\label{factor}
 g(s)=(1-\cos s)(a+b\cos s+c\sin s),\quad
 a=C_0-C_2,\quad b=-2C_2,\quad c=-2S_2.
\end{equation}
For $s\notin2\pi\mathbb Z$ the second factor is nonnegative because
$g(s)\geq0$; continuity gives the same conclusion at the remaining points.
Its minimum is $a-\sqrt{b^2+c^2}$, so $a\geq\sqrt{b^2+c^2}$.
Consequently
\[
 H=\frac12\begin{pmatrix}a&b-ic\\b+ic&a\end{pmatrix}\geq0,
 \qquad T=\begin{pmatrix}1&-1&0\\0&1&-1\end{pmatrix}.
\]
Indeed, the eigenvalues of $H$ are $(a\pm\sqrt{b^2+c^2})/2$.
For $z=e^{is}$ and $w=(1,z,z^2)^T$, one has
$Tw=(1-z)(1,z)^T$ and $|1-z|^2/2=1-\cos s$.
Thus \eqref{factor} equals $w^*(T^*HT/2)w$.
Adding $m e_0e_0^*$, with $e_0=(1,0,0)^T$, gives a positive matrix
$K=m e_0e_0^*+T^*HT/2$ representing $f(s+\alpha)$.
Now set $L=\diag(1,e^{-i\alpha},e^{-2i\alpha})$ and $Q=L^*KL$.
Then $Q\geq0$ and $f(\theta)=w_\theta^*Qw_\theta=q_Q(\theta)$,
where $w_\theta=(1,e^{i\theta},e^{2i\theta})^T
=e^{i\theta}v(\theta)$; the last equality removes a common unit phase.
\end{proof}

\begin{lemma}\label{normalization}
Suppose $Q_a\geq0$ and $\sum_a q_{Q_a}(\theta)=1$ for all $\theta$.
There are $B_a\geq0$ with $q_{B_a}=q_{Q_a}$ such that
\[
 \sum_a B_a=C=\diag(c_0,c_1,c_2),\qquad
 c_j\geq0,\quad c_0+c_1+c_2=1,\quad c_0=c_2.
\]
\end{lemma}
\begin{proof}
Let $J$ reverse the three coordinates, and define
\[
 B_a=\tfrac12(Q_a+JQ_a^TJ).
\]
Transposition and conjugation by $J$ preserve positivity, so $B_a\geq0$.
Since $J\overline{v(\theta)}=v(\theta)$, transposing the scalar quadratic
form shows $q_{JQ^TJ}=q_Q$. Entrywise, $B_{01}=B_{12}$ and
$B_{00}=B_{22}$. These relations pass to $C=\sum_a B_a$.
On the other hand,
\[
 q_C(\theta)=\tr C+(C_{01}+C_{12})e^{i\theta}
 +(C_{10}+C_{21})e^{-i\theta}
 +C_{02}e^{2i\theta}+C_{20}e^{-2i\theta}=1.
\]
Uniqueness of the Fourier coefficients gives $\tr C=1$,
$C_{02}=C_{20}=0$ and $C_{01}+C_{12}=C_{10}+C_{21}=0$.
Together with the entry equalities, these make $C$ diagonal.
Its diagonal is nonnegative because $C\geq0$.
\end{proof}

\section{One state and one measurement}
\begin{proof}[Proof of Theorem~\ref{main}]
Apply Lemmas~\ref{scalar} and~\ref{normalization} to $p$.
Put $D=\diag(\sqrt{c_0},\sqrt{c_1},\sqrt{c_2})$ and let $P$ project onto
the coordinates with $c_j>0$. Let $M$ be diagonal with entry
$c_j^{-1/2}$ when $c_j>0$ and entry zero otherwise.
If $c_j=0$, positivity and $\sum_a(B_a)_{jj}=0$ imply
$(B_a)_{jj}=0$ for every $a$. Each two-coordinate principal minor then
gives $|(B_a)_{jk}|^2\leq(B_a)_{jj}(B_a)_{kk}=0$.
Hence $PB_aP=B_a$, including on the singular boundary.

Writing $r=|\mathcal A|$, define
\begin{equation}\label{effects}
 E_a=MB_aM+\frac1r(I-P).
\end{equation}
Both summands are positive, and
$\sum_aE_a=MCM+I-P=P+I-P=I$.
Furthermore, $DM=MD=P$ and $D(I-P)=0$, so
$DE_aD=B_a$.
Let $\psi=D(1,1,1)^T$ and $\rho=\psi\psi^*$.
Then $\|\psi\|^2=\tr C=1$ and $U_\theta\psi=Dv(\theta)$.
It follows for every outcome and every real angle that
\[
 \tr(U_\theta\rho U_\theta^*E_a)
 =v(\theta)^*DE_aDv(\theta)=q_{B_a}(\theta)=p_a(\theta).
\]
Since $c_0=c_2$, the state has the displayed form with $t=c_0$.
Conversely, a realization \eqref{born} gives nonnegative probabilities
summing to one. Expanding its finite matrix trace produces frequencies
$j-k\in\{-2,-1,0,1,2\}$, even with arbitrary finite multiplicities.
Thus it belongs to $\mathcal R_1^{\mathcal A}$.
\end{proof}

\begin{corollary}\label{postprocessing}
Every stochastic postprocessing of the outcomes is realized with the same
pure state and rotation representation.
\end{corollary}
\begin{proof}
For a stochastic matrix $K(b\mid a)\geq0$ with $\sum_bK(b\mid a)=1$,
the effects $F_b=\sum_aK(b\mid a)E_a$ are positive, sum to $I$, and give
$\sum_a K(b\mid a)p_a(\theta)$ by linearity of the trace.
\end{proof}

\section{Consequences and the dimension boundary}
\begin{corollary}\label{lower}
For every finite nonempty outcome set,
$\mathcal Q_0^{\mathcal A}=\mathcal R_0^{\mathcal A}$ and
$\mathcal Q_{1/2}^{\mathcal A}=\mathcal R_{1/2}^{\mathcal A}$.
These equalities, proved in \cite[\S IV.A, Theorem~5]{AGJLM}, admit
one-dimensional and fixed equatorial-qubit realizations, respectively.
\end{corollary}
\begin{proof}
For constants $p_a=c_a$, take the scalar state $1$ and effects $c_a$.
For $p_a=c_a+u_a\cos\theta+v_a\sin\theta$, positivity gives
$c_a\geq\sqrt{u_a^2+v_a^2}$. The matrices
\[
 E_a=\begin{pmatrix}c_a&u_a-iv_a\\u_a+iv_a&c_a\end{pmatrix}
\]
are therefore positive, and normalization gives $\sum_a E_a=I_2$.
With $\psi=(1,1)^T/\sqrt2$ and
$U_\theta=\diag(e^{-i\theta/2},e^{i\theta/2})$, their Born probabilities
are exactly $p_a$. The reverse inclusions follow by expanding the trace.
\end{proof}

\begin{corollary}\label{sdp}
At spin one, the positive-Gram feasibility conditions for
$\mathcal R_1^{\mathcal A}$ in \cite[Lemma~10 and Appendix~B.2]{AGJLM}
also characterize $\mathcal Q_1^{\mathcal A}$.
\end{corollary}
\begin{proof}
Those conditions ask for $Q_a\geq0$ representing the prescribed Fourier
coefficients, with the polynomials summing to one. Lemma~\ref{normalization}
and \eqref{effects} turn every feasible family into a quantum realization.
Conversely a quantum realization supplies nonnegative normalized
polynomials, hence feasible Gram matrices by Lemma~\ref{scalar}.
Unlike the general fixed-spectrum formulation in \cite{Navascues}, at
spin one no additional diagonal-sum feasibility constraint is needed.
\end{proof}

\begin{corollary}\label{dimension}
Dimension three is necessary for some spin-one families. In particular,
the binary family $p_+=\tfrac12+(\cos\theta+\cos2\theta)/8$,
$p_-=1-p_+$ has no realization in dimension at most two with the
spin-one spectrum.
\end{corollary}
\begin{proof}
Both probabilities lie in $[1/4,3/4]$, so Theorem~\ref{main} applies.
A space of dimension at most two contains at most two distinct frequency
labels. Its trace expansion has at most one nonzero positive frequency
difference, whereas $p_+$ has nonzero coefficients at both one and two.
Uniqueness of Fourier coefficients rules out such a realization.
\end{proof}

\begin{corollary}\label{fixedstate}
There is no fixed density state in any fixed finite spin-one representation
that realizes every binary spin-one family merely by changing the POVM.
In particular, the parameter $t$ in Theorem~\ref{main} cannot be universal.
\end{corollary}
\begin{proof}
Suppose a fixed state $\rho=\sum_k\lambda_k\phi_k\phi_k^*$, with
$\lambda_k>0$ and unit $\phi_k$, realizes
$p_+(\theta)=(1+\cos2\theta)/2$ with an effect $0\leq E\leq I$.
The values one at zero and zero at $\pi/2$ imply
$E\phi_k=\phi_k$ and $EU_{\pi/2}\phi_k=0$ for every $k$.
Hence $\langle\phi_k,U_{\pi/2}\phi_k\rangle=0$.
Decompose $\phi_k$ into the frequency subspaces $-1,0,1$.
The real part of this inner product is $\|(\phi_k)_0\|^2$, so every
middle-frequency component vanishes. Consequently every measurement on
the orbit of $\rho$ has only frequencies zero and $\pm2$, and cannot give
$(1+\cos\theta)/2$. For the balanced pure state, absence of the middle
component specifically forces $t=1/2$.
\end{proof}
This orbit obstruction complements the discussion in \cite[\S IV.B.3]{AGJLM}.
It is not a refutation of an assertion made there. The stronger
no-linear-embedding theorem for the entire spin-one GPT remains the
prior result \cite[Lemma~18]{AGJLM}.

\begin{corollary}\label{boundary}
In dimension four, reversal-transpose averaging alone need not turn a
normalized Gram matrix into a diagonal one.
\end{corollary}
\begin{proof}
Let $w=(1,e^{i\theta},e^{2i\theta},e^{3i\theta})^T$ and set
\[
 C=\tfrac14 I_4+\tfrac1{16}
 \begin{pmatrix}0&1&0&0\\1&0&-2&0\\0&-2&0&1\\0&0&1&0\end{pmatrix}.
\]
The absolute off-diagonal row sums are at most $3/16$; the inequality
$2|\overline{x_i}x_j|\leq|x_i|^2+|x_j|^2$ therefore gives
$x^*Cx\geq\|x\|^2/16$. Thus $C$ is positive. Its off-diagonal sums
at each nonzero frequency vanish, so $w^*Cw=1$. Yet $J_4 C^T J_4=C$
is not diagonal, where $J_4$ reverses the four coordinates.
\end{proof}
The implication tested in Corollary~\ref{boundary} is not asserted in the
source. The higher-spin gaps of \cite[\S IV.C]{AGJLM} remain unchanged.

\section{Coverage and verification}
\begin{longtable}{@{}p{.42\linewidth}p{.53\linewidth}@{}}
\toprule
Source item in \cite{AGJLM} & Coverage\\
\midrule
End of \S IV.B.3, p.~17; repeated in \S VII, p.~27:
$|\mathcal A|\geq3$
 & Answered: yes, by Theorem~\ref{main}.\\
Theorem~1: pure $(2J+1)$-dimensional quantum realization
 & Prior proof cited. Theorem~\ref{main} additionally supplies balanced
 outer populations for every general spin-one family.\\
Theorem~6: binary $\mathcal Q_1=\mathcal R_1$
 & Reproved; prior proof by Aloy et al.\ (2023 preprint; 2024 journal) credited.\\
Lemma~9: positive scalar Gram representation
 & Reproved in degree two in Lemma~\ref{scalar}; prior general
 factorization result cited.\\
\S IV.A: spin zero and spin one-half
 & Reproved; prior proofs by Aloy et al.\ (2023/2024),
 Corollary~\ref{lower}, with the correct lower-spin dimensions.\\
Lemma~10, Appendix~B.2 and the discussion after Lemma~11
 & Not posed; proved consequence at spin one: the rotation-box SDP also characterizes
 quantum boxes (Corollary~\ref{sdp}); compare the general quantum
 characterization of Manos et al.\ (2025/2026).\\
\S IV.B.3: dependence on the qutrit orbit
 & The discussion is respected. Not posed as a universal fixed-state
 claim; obstruction proved in Corollary~\ref{fixedstate}.\\
Lemma~18: no finite quantum embedding of the GPT
 & Prior proof cited; compatible with a family-dependent state.\\
\S IV.B.1 and Lemmas~16--17: facial structure and distinguishability
 & Prior results not covered; no change to the source's GPT distinctions.\\
Dimension optimality; stochastic postprocessing
 & Not posed; proved in Corollaries~\ref{dimension} and~\ref{postprocessing}.\\
Reversal normalization at degree three
 & Not asserted in the source; the diagonalization implication fails
 for the matrix in Corollary~\ref{boundary}.\\
\S IV.C: higher-spin gaps; \S IV.D--F: approximation, two-angle
 restrictions and classicality; \S V: Bell correlations;
 \S VI: representation-theoretic reformulations
 & Prior results not covered by this construction.\\
\bottomrule
\end{longtable}

The universal statements and the scalar Gram construction have written
proofs above. The pinned Lean declaration
\begin{quote}\small\ttfamily
D5.S3.Quantum.Measurement.RotationBoxes.\\
SpinOneRotationBoxes.result
\end{quote}
states the equivalence of
\texttt{QuantumSpinOne} and \texttt{RotationSpinOne} for every finite
outcome type of cardinality at least three; its definitions include finite
frequency multiplicities, a density state and a single normalized positive
measurement. The exact source is
\href{https://github.com/the-omega-institute/trureturing/blob/4b14c77599edfd8a52c0a7926d038f4121ca7db1/D5/S3/Quantum/Measurement/RotationBoxes/SpinOneRotationBoxes.lean}{\texttt{SpinOneRotationBoxes.lean} at \texttt{4b14c77599ed}}.
The stronger displayed pure-state and balanced-population form is supplied
by the written construction, rather than by the statement of that declaration.

Within the same file, \texttt{simultaneous\_\allowbreak diagonal\_\allowbreak normalization}
states preservation of the Gram polynomials, positivity, a diagonal
trace-one sum and equality of its outer diagonal entries;
\texttt{gram\_\allowbreak family\_\allowbreak quantum} states the passage from positive
diagonal-sum matrices to a quantum realization. Both are auxiliary
 declarations. The one-outcome extension and Corollaries~\ref{lower}--\ref{boundary}
have the written proofs given here.

The accompanying programs check the scalar factorization,
reversal-transpose and support-normalization identities, the source's
binary example, a specified finite catalog of three- and four-outcome
Gaussian-integer Gram families, and the displayed corollary formulas.
These finite and symbolic checks support the algebraic steps; the
universal conclusions follow from the written proofs, not enumeration.

\begin{thebibliography}{9}
\bibitem{AGJLM}
A. Aloy, T. D. Galley, C. L. Jones, S. L. Ludescher and M. P. M\"uller,
\emph{Spin-bounded correlations: rotation boxes within and beyond quantum theory},
Communications in Mathematical Physics \textbf{405}, 292 (2024).
\href{https://doi.org/10.1007/s00220-024-05123-2}{doi:10.1007/s00220-024-05123-2};
\href{https://arxiv.org/abs/2312.09278v2}{arXiv:2312.09278v2}.
First preprint: \href{https://arxiv.org/abs/2312.09278v1}{v1}, December 14, 2023.

\bibitem{Navascues}
K. Manos, M. Weilenmann and M. Navascu\'es,
\emph{Timed demolition measurements},
\href{https://arxiv.org/abs/2507.06912v2}{arXiv:2507.06912v2}
(2025; revised March 4, 2026), \S III.B, Eqs.~(44)--(47).
\end{thebibliography}
\end{document}
